\documentclass[10pt,a4paper]{amsart}
\usepackage[margin=19mm]{geometry}
\usepackage{amsmath,amssymb,mathtools}
\usepackage[scr=rsfs,cal=euler]{mathalfa}
\usepackage{xfrac}
\usepackage[unicode,hidelinks,bookmarks=false]{hyperref}
\newcommand{\ie}{i.e.\ }
\newcommand{\cf}{cf.\ }
\newcommand{\resp}{respectively\ }
\newcommand{\Subsection}{\subsection}
\providecommand{\nobreakdash}{\nobreak\hbox{-}}
\setlength{\emergencystretch}{2em}
\multlinegap0pt
\usepackage{amssymb}
\usepackage{xcolor}
\usepackage{algorithm}\usepackage{algpseudocode}
\usepackage[shortlabels]{enumitem}
\newtheorem{definition}{Definition}
\newtheorem{lemma}{Lemma}
\newtheorem{theorem}{Theorem}
\title{Constrained Nonnegative Gram Feasibility is $\exists\mathbb{R}$-Complete}
\author{Angshul Majumdar, IIIT Delhi, India}
\date{Source: arXiv:2603.19976v1 (CC BY 4.0)}
\begin{document}
\maketitle

\noindent\emph{Source and licence.} Constrained Nonnegative Gram Feasibility is $\exists\mathbb{R}$-Complete. Version arXiv:2603.19976v1 (CC BY 4.0), distributed by arXiv under the Creative Commons Attribution 4.0 International licence, http://creativecommons.org/licenses/by/4.0/, as recorded in the arXiv licence metadata for this exact version and shown on the abstract page https://arxiv.org/abs/2603.19976v1. Full licence notice: \texttt{licenses/arxiv-2603.19976-CC-BY-4.0.txt}.\par\smallskip

\noindent\textit{Editorial scope.} Counted unit: the article's theorem thm:hardness-fixed-r (source0.tex lines 1208--1212). The article's complete original proof of it is delivered with it (source0.tex lines 1214--1369, one \texttt{proof} environment of 156 lines). No mathematics is written, repaired or re-derived: the statement and the proof are the article's own lines, in the article's own notation, and no retained line is substituted or re-flowed.  Retained uncounted as the source-local context the proof needs: lines 246--260; lines 1139--1159.  Nothing was omitted from any retained span, so the package carries no omission record.  No retained line contains a citation, so the wrapper carries no bibliography and no bibliography entry.  No retained line contains a figure environment, an image-inclusion command or drawing code, so no image is delivered and the package declares no asset dependency.  Editorial formatting only, in addition: an AMS \texttt{amsart} preamble that restates the article's own declarations and macros used by the retained lines, namely the environments \texttt{definition}, \texttt{lemma}, \texttt{theorem} on separate counters, together with the standard packages those lines need.  The retained environments print in this document as Definition 1, Lemma 1, Theorem 1 in order of appearance, whereas the article prints the same items with the numbering of its own full document.  The complete native source of the article as distributed in the arXiv e-print for this exact version is bundled unchanged as \texttt{sources/196821/source0.tex}.
\begin{definition}[ETR-AMI]
\label{def:etrami}
An instance of \textsc{ETR-AMI} consists of variables
\[
x_1,\dots,x_n \in \mathbb{R}_{\ge 0}
\]
together with constraints of the following three forms:
\[
x_i = 1,\qquad
x_i = x_j + x_k,\qquad
x_i = x_j x_k.
\]
The decision problem asks whether there exists a nonnegative real assignment
satisfying all constraints.
\end{definition}

\begin{lemma}[Anchor frame rigidity]
\label{lem:anchor-frame-rigidity}
Let \(r\ge 2\), and suppose a symmetric matrix \(W\) admits a rank-$r$
nonnegative Gram realization with distinguished indices
\[
e_1,\dots,e_r
\]
satisfying
\[
W_{e_a e_b}=\delta_{ab}
\qquad \text{for all } a,b\in\{1,\dots,r\},
\]
where \(\delta_{ab}\) denotes the Kronecker delta. Then, up to a common
permutation of the \(r\) coordinates,
\[
h_{e_a}=e_a^{(r)}
\qquad \text{for all } a=1,\dots,r,
\]
where \(e_a^{(r)}\) is the \(a\)-th standard basis vector of
\(\mathbb{R}^r\).
\end{lemma}

\begin{theorem}
\label{thm:hardness-fixed-r}
For every fixed integer \(r\ge 2\), constrained rank-$r$ nonnegative Gram
feasibility is \(\exists\mathbb{R}\)-hard.
\end{theorem}

\begin{proof}
Fix \(r\ge 2\). We reduce from \textsc{ETR-AMI}
(Definition~\ref{def:etrami}).

Let
\[
\mathcal{I}
\]
be an instance of \textsc{ETR-AMI} with variables
\[
x_1,\dots,x_n\in\mathbb{R}_{\ge 0}
\]
and constraints of the forms
\[
x_i=1,\qquad x_i=x_j+x_k,\qquad x_i=x_jx_k.
\]

We construct an instance of constrained rank-$r$ nonnegative Gram feasibility
as follows.

\paragraph{Anchor rows.}
Introduce \(r\) distinguished indices
\[
e_1,\dots,e_r
\]
and specify the Gram entries
\[
\widehat{W}_{e_a e_b}=\delta_{ab}
\qquad \text{for all } a,b\in\{1,\dots,r\}.
\]
By Lemma~\ref{lem:anchor-frame-rigidity}, every feasible realization may be
normalized so that
\[
h_{e_a}=e_a^{(r)}
\qquad \text{for all } a=1,\dots,r.
\]

\paragraph{Variable rows.}
For each arithmetic variable \(x_i\), introduce one index \(u_i\).
Impose the specified Gram entry
\[
\widehat{W}_{e_2 u_i}=1,
\]
and, for every \(t=3,\dots,r\), impose the specified Gram entry
\[
\widehat{W}_{e_t u_i}=0.
\]

We claim that in every feasible realization,
\[
h_{u_i}=(x_i,1,0,\dots,0)
\]
for a uniquely determined scalar \(x_i\in\mathbb{R}_{\ge 0}\).
Indeed, write
\[
h_{u_i}=(a_{i,1},a_{i,2},\dots,a_{i,r})\in\mathbb{R}_+^r.
\]
Since \(h_{e_2}=e_2^{(r)}\), the constraint \(\widehat{W}_{e_2 u_i}=1\) gives
\[
a_{i,2}=1.
\]
Likewise, for each \(t=3,\dots,r\), the constraint
\(\widehat{W}_{e_t u_i}=0\) gives
\[
a_{i,t}=0.
\]
Setting \(x_i:=a_{i,1}\), we obtain
\[
h_{u_i}=(x_i,1,0,\dots,0),
\]
as claimed.

\paragraph{Constant constraints.}
For every source constraint \(x_i=1\), impose the affine side constraint
\[
W_{e_1 u_i}=1.
\]

\paragraph{Addition constraints.}
For every source constraint \(x_i=x_j+x_k\), impose the affine side constraint
\[
W_{e_1 u_j}+W_{e_1 u_k}-W_{e_1 u_i}=0.
\]

\paragraph{Multiplication constraints.}
For every source constraint \(x_i=x_jx_k\), impose the affine side constraint
\[
W_{u_j u_k}-W_{e_1 u_i}-1=0.
\]

We now verify correctness.

\paragraph{Soundness.}
Suppose \(\mathcal{I}\) has a satisfying assignment
\[
x_1,\dots,x_n\in\mathbb{R}_{\ge 0}.
\]
Define
\[
h_{e_a}=e_a^{(r)}
\qquad \text{for } a=1,\dots,r,
\]
and
\[
h_{u_i}=(x_i,1,0,\dots,0)
\qquad \text{for } i=1,\dots,n.
\]
These vectors lie in \(\mathbb{R}_+^r\). The anchor constraints hold by
construction. The variable-row constraints hold because
\[
\langle e_2^{(r)},(x_i,1,0,\dots,0)\rangle = 1
\]
and
\[
\langle e_t^{(r)},(x_i,1,0,\dots,0)\rangle = 0
\qquad \text{for } t=3,\dots,r.
\]
Moreover,
\[
W_{e_1 u_i}=x_i,
\]
and
\[
W_{u_j u_k}=x_jx_k+1.
\]
Hence the constant, addition, and multiplication constraints are satisfied
exactly as in the rank-$2$ construction.

\paragraph{Completeness.}
Conversely, suppose the constructed instance has a feasible realization.
By Lemma~\ref{lem:anchor-frame-rigidity}, we may normalize so that
\[
h_{e_a}=e_a^{(r)}
\qquad \text{for all } a=1,\dots,r.
\]
The variable-row constraints then force
\[
h_{u_i}=(x_i,1,0,\dots,0)
\]
for some uniquely determined \(x_i\in\mathbb{R}_{\ge 0}\), as shown above.
Consequently,
\[
W_{e_1 u_i}=x_i,
\qquad
W_{u_j u_k}=x_jx_k+1.
\]
The affine side constraints therefore enforce exactly the source arithmetic
constraints \(x_i=1\), \(x_i=x_j+x_k\), and \(x_i=x_jx_k\). Thus the extracted
tuple \((x_1,\dots,x_n)\) is a satisfying assignment of \(\mathcal{I}\).

Therefore the reduction is correct. Since \(r\) is fixed, the construction
uses \(n+r\) indices and \(O(n+m+r^2)\) specified entries and affine
constraints, where \(m\) is the number of source constraints. Hence the
construction is polynomial-time. It follows that constrained rank-$r$
nonnegative Gram feasibility is \(\exists\mathbb{R}\)-hard.
\end{proof}

\end{document}
